\documentclass[12pt]{article}
\usepackage[utf8]{inputenc}
\usepackage[T1]{fontenc}
\usepackage{hyperref}
\title{Yirrk-Mel Language Resources for K-12 Curriculum Development}
\author{Compiled by Claude Sonnet 5.5 (claude-sonnet-5-5) for Andr\'as Kornai}
\date{}
\begin{document}
\maketitle

\section{Introduction}
Yirrk-Mel (ISO 639-3 yrm) is also called Yirrk-Thangalkl. Search results describe it as a dialect of Yir-Yoront, a Paman language of the south-western Cape York Peninsula, Queensland. Sources list it as extinct or dormant. Lauriston Sharp wrote the name as ``Yir Mel''. Yir-Yoront itself had only about 15 speakers in 1991. I found no school materials, primers or community programmes for Yirrk-Mel in the searches. The few resources below are general reference pages, and they mostly cover the closely related Yir-Yoront.

\section{Primary Educational Resources}
No primers, textbooks or school programmes for Yirrk-Mel were found.

\section{Linguistic and Reference Materials}
\begin{itemize}
\item Wikipedia, Yirrk-Thangalkl dialect: \url{https://en.wikipedia.org/wiki/Yirrk-Thangalkl_dialect}
\item Wikipedia, Yir-Yoront language (the closely related language): \url{https://en.wikipedia.org/wiki/Yir-Yoront_language}
\item AIATSIS Austlang entry: \url{https://collection.aiatsis.gov.au/austlang/language/y214}
\end{itemize}
These are descriptive overviews. I did not open them to check whether they list grammars, wordlists or archival recordings.

\section{Religious and Community Resources}
None were found in the searches.

\section{Digital and Audio Resources}
None were found. The AIATSIS collection may hold archival material on Yir-Yoront. I did not verify this.

\section{Curriculum Development Context}
The language is dormant and no community of first-language speakers remains. Any curriculum would be a revitalisation project. It would have to rely on archival documentation of Yir-Yoront and Yirrk-Thangalkl, and it would need the involvement of the descendant community (Kowanyama area) and AIATSIS.
\end{document}